\begin{tikzpicture}[
    node distance=0.7cm and 1.0cm,
    every node/.style={font=\scriptsize},
    edge/.style={rectangle, rounded corners, draw=black!70, very thick, fill=green!18, text width=2.0cm, align=center, minimum height=0.8cm},
    cloud/.style={rectangle, rounded corners, draw=black!70, very thick, fill=orange!22, text width=2.2cm, align=center, minimum height=0.8cm},
    db/.style={cylinder, draw=black!70, fill=blue!12, aspect=0.3, shape border rotate=90, minimum height=1.2cm, minimum width=2.0cm, font=\scriptsize\bfseries},
    arrow/.style={->, >=stealth, thick}
]

% --- Edge ---
\node[edge] (ad)   {AD8232\\analog};
\node[edge, right=of ad] (adc) {ADC1\\DMA};
\node[edge, right=of adc] (down) {downsample\\500$\to$100\,Hz};
\node[edge, right=of down] (ble) {BLE 0xFF03\\(int16, 10/pkt)};

% --- Gateway de testing ---
\node[edge, right=1.2cm of ble, fill=cyan!18] (app) {Cliente PC\\\texttt{ble\_logger.py}};
\node[db, right=1.4cm of app] (fs) {Firestore};

% --- Pan-Tompkins pipeline ---
\node[cloud, below=0.8cm of fs] (fn) {process\_ecg};
\node[cloud, below=0.7cm of fn] (bp)   {Bandpass\\5--15\,Hz};
\node[cloud, left=of bp] (deriv) {Derivativo\\5-puntos};
\node[cloud, left=of deriv] (sq) {Cuadrado};
\node[cloud, left=of sq] (mw) {MWI 150\,ms};
\node[cloud, left=of mw] (peaks) {Detección\\R-peaks};
\node[cloud, left=of peaks] (bpm) {BPM, HRV\\(SDNN)};

% --- Connections edge -> cloud ---
\draw[arrow] (ad) -- (adc);
\draw[arrow] (adc) -- (down);
\draw[arrow] (down) -- (ble);
\draw[arrow] (ble) -- (app);
\draw[arrow] (app) -- node[above, font=\scriptsize, inner sep=2pt]{Wi-Fi} (fs);
\draw[arrow] (fs)  -- (fn);
\draw[arrow] (fn)  -- (bp);
\draw[arrow] (bp)  -- (deriv);
\draw[arrow] (deriv)  -- (sq);
\draw[arrow] (sq)  -- (mw);
\draw[arrow] (mw)  -- (peaks);
\draw[arrow] (peaks) -- (bpm);

% --- Solución de la línea cruzada (Loop inferior de retorno) ---
% Sale por abajo de BPM, recorre todo hacia la derecha por debajo del pipeline, y sube a Firestore.
\draw[arrow, rounded corners=6pt] (bpm.south) -- ++(0,-0.6) coordinate(ret1) -| ([xshift=0.8cm]bp.east) coordinate(ret2) |- (fs.0);

% --- Section bands ---
\begin{scope}[on background layer]
    % Coordenadas falsas para dar margen superior a los títulos de las cajas
    \path (ad.north) ++(0, 0.45) coordinate (dummyEdge);
    \path (fs.north) ++(0, 0.55) coordinate (dummyCloud);

    \node[rectangle, draw=green!50!black, dashed, inner sep=0.35cm,
          fit=(ad) (ble) (dummyEdge), fill=green!4] (edgeBox) {};
    \node[anchor=north west, font=\scriptsize\bfseries, green!50!black] 
          at ([xshift=4pt, yshift=-4pt]edgeBox.north west) {Edge (ESP32)};

    \node[rectangle, draw=cyan!60!black, dashed, inner sep=0.35cm,
          fit=(app), fill=cyan!4] (appBox) {};
    \node[anchor=north, font=\scriptsize\bfseries, cyan!60!black] 
          at ([yshift=-4pt]appBox.north) {Gateway};

    % Al agregar las coordenadas (ret1) y (ret2) al 'fit', la caja crece para envolver la nueva línea inferior.
    \node[rectangle, draw=orange!70!black, dashed, inner sep=0.35cm,
          fit=(fn) (bp) (bpm) (fs) (ret1) (ret2) (dummyCloud), fill=orange!4] (cloudBox) {};
    \node[anchor=north east, font=\scriptsize\bfseries, orange!70!black] 
          at ([xshift=-4pt, yshift=-4pt]cloudBox.north east) {Firebase Cloud Functions (Python/NumPy)};
\end{scope}

\end{tikzpicture}